% !TEX root = ../main.tex
\begin{ouverture}[champuniforme]{15}{Mouvement dans un champ uniforme}{Partie 2 · Mouvement et interactions}
{Champ de pesanteur\enspace·\enspace Équations horaires et trajectoire\enspace·\enspace Champ électrique d'un condensateur\enspace·\enspace Énergie}
\objectif{Montrer que le mouvement dans un champ uniforme est plan.}
\objectif{Établir et exploiter les équations horaires du mouvement et l'équation de la trajectoire.}
\objectif{Discuter de l'influence des grandeurs physiques sur le champ électrique d'un condensateur plan.}
\objectif{Décrire le principe d'un accélérateur linéaire de particules chargées.}
\objectif{Exploiter la conservation de l'énergie mécanique ou le théorème de l'énergie cinétique.}
\prerequis{Deuxième loi de Newton (chapitre~14) ; vecteurs vitesse et accélération (chapitre~13).}
\prerequis{Énergies cinétique, potentielle et mécanique ; travail d'une force (classe de première).}
\prerequis{Primitives des fonctions $t\mapsto a$ et $t\mapsto a\,t+b$ ; équation du second degré.}
\end{ouverture}
\essentiel
\essSection{1}{Mouvement dans le champ de pesanteur}
\begin{propriete}[Accélération dans un champ de pesanteur uniforme]
Un projectile soumis à son seul poids (frottements négligés) vérifie, dans le référentiel terrestre supposé galiléen,
$m\,\vec{a}=m\,\vec{g}$, soit $\vec{a}=\vec{g}$ : son accélération est constante, verticale, dirigée vers le bas, et ne
dépend pas de sa masse. Le mouvement est \textbf{plan} : il reste dans le plan vertical qui contient le vecteur vitesse
initiale $\vec{v}_0$.
\end{propriete}
\begin{savoirfaire}[Établir les équations horaires]
Choisir un repère $(O\,;\vec{\imath},\vec{k})$ dans le plan du mouvement, $\vec{k}$ vertical vers le haut, $O$ au sol à la
verticale du point de lancement, situé à la hauteur $h$. Avec $\vec{v}_0$ incliné d'un angle $\alpha$ :
$a_x=0$ et $a_z=-g$ ; on cherche des primitives en tenant compte des conditions initiales,
\[v_x=v_0\cos\alpha,\quad v_z=-g\,t+v_0\sin\alpha\qquad\text{puis}\qquad x=v_0\cos\alpha\;t,\quad z=-\frac12\,g\,t^2+v_0\sin\alpha\;t+h.\]
En éliminant $t=\dfrac{x}{v_0\cos\alpha}$, on obtient l'équation de la trajectoire, un arc de parabole :
\[z=-\frac{g}{2\,v_0^2\cos^2\alpha}\,x^2+x\tan\alpha+h.\]
\end{savoirfaire}
\schema{15-parabole}
\begin{exemple}
Une boule de pétanque est lancée d'une hauteur $h=\qty{1,0}{m}$ avec $v_0=\qty{7,0}{m.s^{-1}}$ et $\alpha=40^\circ$. Au
sommet, $v_z=0$ : $t_\mathrm{S}=\dfrac{v_0\sin\alpha}{g}=\qty{0,46}{s}$ et la \textbf{flèche} vaut $z_\mathrm{S}=\qty{2,0}{m}$.
La boule touche le sol quand $z=0$ : la résolution de l'équation du second degré donne $t=\qty{1,10}{s}$, d'où la
\textbf{portée} $x=v_0\cos\alpha\;t=\qty{5,9}{m}$.
\end{exemple}
\essSection{2}{Mouvement dans un champ électrique uniforme}
\begin{propriete}[Champ électrique d'un condensateur plan]
Entre deux plaques parallèles distantes de $d$, soumises à la tension $U$, le champ électrique est uniforme :
perpendiculaire aux plaques, dirigé de la plaque positive vers la plaque négative, de valeur
\[E=\frac{U}{d}\qquad(E\text{ en }\unit{V.m^{-1}}).\]
Il est proportionnel à la tension et inversement proportionnel à la distance entre les plaques. Une particule de
charge $q$ y subit la force $\vec{F}=q\,\vec{E}$.
\end{propriete}
\schema{15-condensateur}
\begin{propriete}[Mouvement d'une particule chargée]
Pour une particule chargée (électron, proton, ion), le poids est négligeable devant la force électrique : pour un
électron dans un champ de \qty{1e4}{V.m^{-1}}, $F=e\,E=\qty{1,6e-15}{N}$, alors que $P=m\,g=\qty{8,9e-30}{N}$. La
deuxième loi donne $\vec{a}=\dfrac{q}{m}\,\vec{E}$, constant : comme dans le champ de pesanteur, le mouvement est plan
et la trajectoire est rectiligne ou parabolique selon la direction de la vitesse initiale.
\end{propriete}
\begin{definition}[Accélérateur linéaire]
Dans un \textbf{accélérateur linéaire}, des particules chargées traversent successivement des intervalles entre
électrodes, soumis à une tension $U$ : à chaque traversée, leur énergie cinétique augmente de $|q|\,U$. Entre deux
intervalles, elles circulent dans des tubes métalliques où le champ est nul ; la tension alternative s'inverse pendant
ce trajet, si bien que les particules sont toujours accélérées.
\end{definition}
\essSection{3}{Aspects énergétiques}
\begin{propriete}[Théorème de l'énergie cinétique et énergie mécanique]
Entre deux points A et B, la variation de l'énergie cinétique est égale à la somme des travaux des forces :
$\Delta E_\mathrm{c}=E_\mathrm{c}(B)-E_\mathrm{c}(A)=\sum W_{AB}(\vec{F})$, avec $W_{AB}(\vec{P})=m\,g\,(z_\mathrm{A}-z_\mathrm{B})$ et
$W_{AB}(q\vec{E})=q\,U_{AB}$. Si seul le poids travaille, l'énergie mécanique $E_\mathrm{m}=\frac12\,m\,v^2+m\,g\,z$ se
\textbf{conserve}.
\end{propriete}
\begin{exemple}
Un électron, initialement immobile, est accéléré sous une tension $U=\qty{2,0}{kV}$ : $E_\mathrm{c}=e\,U=\qty{3,2e-16}{J}$,
d'où $v=\sqrt{\dfrac{2\,e\,U}{m}}=\qty{2,7e7}{m.s^{-1}}$, près du dixième de la vitesse de la lumière.
\end{exemple}
\begin{python}
m, g, tau = 0.058, 9.81, 0.040                         # kg, N/kg, s
x = [0.000, 0.120, 0.240, 0.360, 0.480, 0.600, 0.720]  # m
z = [1.000, 1.152, 1.289, 1.409, 1.514, 1.604, 1.677]  # m
for i in range(1, 6):
    vx = (x[i+1] - x[i-1]) / (2 * tau)
    vz = (z[i+1] - z[i-1]) / (2 * tau)
    Ec = 0.5 * m * (vx**2 + vz**2)
    Ep = m * g * z[i]
    print(round(Ec, 2), round(Ep, 2), round(Ec + Ep, 2))
\end{python}
\begin{console}
0.64 0.66 1.29
0.56 0.73 1.29
0.49 0.8 1.29
0.43 0.86 1.29
0.38 0.91 1.29
\end{console}
\noindent\capacite{Capacité numérique}Les listes contiennent les coordonnées d'une balle de tennis de 58~g pointées
sur une vidéo. Le programme calcule à chaque position l'énergie cinétique, l'énergie potentielle de pesanteur et leur
somme : pendant la montée, $E_\mathrm{c}$ diminue, $E_\mathrm{pp}$ augmente et $E_\mathrm{m}$ reste égale à
\qty{1,29}{J}. Les frottements sont bien négligeables.
\begin{experience}[Étudier un lancer par vidéo]
On filme un lancer de balle devant une règle étalon, puis on pointe les positions avec un logiciel ou un programme. On
modélise $x(t)$ par une fonction affine et $z(t)$ par un polynôme du second degré : le coefficient de $t^2$ vaut
$-g/2$, et les coefficients du premier degré donnent $v_0\cos\alpha$ et $v_0\sin\alpha$. Le tracé des énergies vérifie la
conservation de l'énergie mécanique.
\end{experience}
\begin{rediger}
Pour établir les équations horaires, citer la loi utilisée, écrire les conditions initiales, puis primitiver une
coordonnée à la fois. Pour une question de portée ou de hauteur, préciser la condition qui la caractérise : « au sommet,
$v_z=0$ » ; « au sol, $z=0$ ».
\end{rediger}
\exercices

\rubrique{Découverte}
\begin{blocLivre}{ex-15-1}
\begin{exercice}[1]{Une chute libre verticale}
Une bille est lâchée sans vitesse initiale d'une hauteur de 20~m. On néglige les frottements et on prend un axe $(Oz)$
vertical, orienté vers le bas, d'origine au point de lâcher.
\begin{donnees}
$g=\qty{9,81}{m.s^{-2}}$.
\end{donnees}
\begin{questions}
\item Appliquer la deuxième loi de Newton et en déduire l'accélération de la bille.
\item Établir, par primitives successives, les équations horaires $v_z(t)$ et $z(t)$.
\item Calculer la durée de la chute et la vitesse de la bille quand elle touche le sol.
\item Retrouver cette vitesse par la conservation de l'énergie mécanique.
\end{questions}
\end{exercice}
\end{blocLivre}
\begin{blocLivre}{ex-15-2}
\begin{exercice}[1]{Le champ d'un condensateur plan}
Deux plaques métalliques parallèles, distantes de $d=\qty{2,0}{cm}$, sont soumises à une tension $U=\qty{500}{V}$.
\begin{donnees}
$e=\qty{1,60e-19}{C}$ ; masse de l'électron : $m=\qty{9,11e-31}{kg}$ ; $g=\qty{9,81}{N.kg^{-1}}$.
\end{donnees}
\begin{questions}
\item Caractériser le champ électrique entre les plaques : direction, sens et valeur.
\item Comment varie la valeur du champ si l'on double la tension ? si l'on double la distance entre les plaques ?
\item Calculer la force électrique subie par un électron entre les plaques, puis son poids. Conclure.
\end{questions}
\end{exercice}
\end{blocLivre}
\begin{blocLivre}{ex-15-3}
\begin{exercice}[1]{La bille qui quitte la table}
Une bille roule sur une table horizontale de hauteur \qty{0,80}{m} et la quitte avec une vitesse horizontale de
\qty{2,5}{m.s^{-1}}. On néglige les frottements de l'air.
\begin{questions}
\item Établir les équations horaires du mouvement de la bille, l'origine étant prise au sol, à la verticale du bord de la
table.
\item Calculer la durée de la chute.
\item À quelle distance du pied de la table la bille touche-t-elle le sol ?
\item Calculer la valeur de sa vitesse à l'arrivée au sol.
\end{questions}
\end{exercice}
\end{blocLivre}
\begin{blocLivre}{ex-15-4}
\begin{exercice}[1]{Les énergies d'une balle}
Une balle de masse $m=\qty{0,20}{kg}$ est lancée verticalement vers le haut, depuis une hauteur de \qty{1,5}{m}, à la
vitesse de \qty{8,0}{m.s^{-1}}. On néglige les frottements ; l'énergie potentielle de pesanteur est nulle au sol.
\begin{questions}
\item Calculer l'énergie cinétique, l'énergie potentielle et l'énergie mécanique de la balle au départ.
\item Pourquoi l'énergie mécanique se conserve-t-elle ?
\item En déduire l'altitude maximale atteinte par la balle.
\item Avec quelle vitesse la balle touche-t-elle le sol ?
\end{questions}
\end{exercice}
\end{blocLivre}
\begin{blocLivre}{ex-15-5}
\begin{exercice}[1]{Accélérer un électron}
Dans un canon à électrons, des électrons émis sans vitesse sont accélérés entre deux électrodes soumises à une tension
$U=\qty{1,0}{kV}$.
\begin{donnees}
$e=\qty{1,60e-19}{C}$ ; $m=\qty{9,11e-31}{kg}$.
\end{donnees}
\begin{questions}
\item Exprimer le travail de la force électrique entre les deux électrodes.
\item Appliquer le théorème de l'énergie cinétique et calculer l'énergie cinétique finale.
\item En déduire la vitesse des électrons à la sortie.
\item Comment faudrait-il modifier la tension pour doubler cette vitesse ?
\end{questions}
\end{exercice}
\end{blocLivre}
\rubrique{Entraînement}
\begin{blocLivre}{ex-15-6}
\begin{exercice}[2]{Le lancer du poids}
Un athlète lance le poids depuis une hauteur $h=\qty{2,10}{m}$, avec une vitesse $v_0=\qty{13,5}{m.s^{-1}}$ inclinée de
$\alpha=42^\circ$ sur l'horizontale. On néglige les frottements de l'air.
\begin{questions}
\item Justifier que le mouvement du centre du poids est plan.
\item Établir les équations horaires, puis l'équation de la trajectoire.
\item Calculer la flèche, c'est-à-dire l'altitude maximale.
\item Calculer la longueur du jet, mesurée au sol depuis la verticale du point de lâcher.
\end{questions}
\end{exercice}
\end{blocLivre}
\begin{blocLivre}{ex-15-7}
\begin{exercice}[2]{Réussir son tir}
Un basketteur tire depuis une hauteur de \qty{2,20}{m}, avec un angle $\alpha=52^\circ$. Le centre de l'anneau se trouve à
la distance horizontale \qty{4,20}{m} du point de tir, à \qty{3,05}{m} du sol.
\begin{questions}
\item Rappeler l'équation de la trajectoire du centre du ballon.
\item Montrer que la vitesse initiale nécessaire vérifie
$v_0^2=\dfrac{g\,x^2}{2\cos^2\alpha\,(x\tan\alpha+h-z)}$, avec $x=\qty{4,20}{m}$ et $z=\qty{3,05}{m}$.
\item Calculer $v_0$.
\item Le joueur lance à \qty{7,4}{m.s^{-1}} : le ballon passe-t-il au-dessus ou au-dessous du centre de l'anneau ?
\end{questions}
\end{exercice}
\end{blocLivre}
\begin{blocLivre}{ex-15-8}
\begin{exercice}[2]{Dévier un faisceau d'électrons}
Des électrons pénètrent en O, avec une vitesse horizontale $v_0=\qty{2,0e7}{m.s^{-1}}$, entre deux plaques horizontales de
longueur $L=\qty{5,0}{cm}$, distantes de $d=\qty{2,0}{cm}$, soumises à la tension $U=\qty{200}{V}$.
\schema{15-ex08}
\begin{donnees}
$e=\qty{1,60e-19}{C}$ ; $m=\qty{9,11e-31}{kg}$ ; le poids des électrons est négligeable.
\end{donnees}
\begin{questions}
\item Calculer la valeur du champ électrique et préciser le sens de la force subie par les électrons.
\item Établir les équations horaires du mouvement entre les plaques, puis l'équation de la trajectoire.
\item Calculer la déviation verticale à la sortie des plaques. Les électrons touchent-ils une plaque ?
\item Calculer l'angle que fait leur vitesse avec l'horizontale à la sortie.
\end{questions}
\end{exercice}
\end{blocLivre}
\begin{blocLivre}{ex-15-9}
\begin{exercice}[2]{Tracer l'évolution des énergies}
On reprend le programme de l'essentiel : il calcule les énergies d'une balle de tennis de 58~g à partir d'un pointage
vidéo.
\begin{questions}
\item \capacite{Capacité numérique}Expliquer comment le programme calcule la vitesse de la balle en chaque point.
\item Quelle est l'origine de l'énergie potentielle choisie ? Pourquoi peut-on ignorer les positions extrêmes de la
liste ?
\item Interpréter l'évolution des trois grandeurs affichées.
\item Écrire les lignes à ajouter pour tracer $E_\mathrm{c}$, $E_\mathrm{pp}$ et $E_\mathrm{m}$ en fonction du temps.
\item En présence de frottements notables, comment évoluerait la troisième colonne ?
\end{questions}
\end{exercice}
\end{blocLivre}
\begin{blocLivre}{ex-15-10}
\begin{exercice}[2]{Le toboggan aquatique}
Un enfant de 30~kg se laisse glisser, sans vitesse initiale, du haut d'un toboggan de \qty{8,0}{m} de dénivelé et de
\qty{30}{m} de long. Il arrive en bas à \qty{9,0}{m.s^{-1}}.
\begin{questions}
\item Sans frottement, quelle serait sa vitesse en bas ? Justifier.
\item Comparer à la vitesse mesurée. Que conclure ?
\item Appliquer le théorème de l'énergie cinétique pour calculer le travail des forces de frottement.
\item En supposant la force de frottement constante et parallèle à la piste, calculer sa valeur.
\end{questions}
\end{exercice}
\end{blocLivre}
\begin{blocLivre}{ex-15-11}
\begin{exercice}[2]{L'accélérateur linéaire}
Un accélérateur linéaire comporte une succession de tubes métalliques alignés ; entre deux tubes, les protons traversent
un intervalle soumis à une tension de \qty{50}{kV}, qui les accélère à chaque passage. La tension alternative a pour
fréquence $f=\qty{10}{MHz}$.
\schema{15-ex11}
\begin{donnees}
$e=\qty{1,60e-19}{C}$ ; masse du proton : $m_\mathrm{p}=\qty{1,67e-27}{kg}$ ; $\qty{1}{eV}=\qty{1,60e-19}{J}$.
\end{donnees}
\begin{questions}
\item Quelle énergie cinétique un proton gagne-t-il à chaque intervalle ? L'exprimer en joules et en électronvolts.
\item Les protons entrent avec une vitesse négligeable. Calculer leur vitesse après 1 intervalle, puis après 20.
\item À l'intérieur d'un tube, le champ est nul. Quelle est la nature du mouvement ?
\item Pour être accélérés au bon moment, les protons doivent traverser chaque tube en une demi-période de la tension.
Calculer la longueur du premier tube et celle du vingtième. Pourquoi les tubes sont-ils de plus en plus longs ?
\end{questions}
\end{exercice}
\end{blocLivre}
\rubrique{Objectif bac}
\begin{blocLivre}{ex-15-12}
\begin{exercice}[3]{Le service au tennis}
\begin{docu}{1}{Le terrain}
Le serveur frappe la balle à \qty{2,80}{m} de hauteur, sur sa ligne de fond. Le filet, haut de \qty{0,914}{m}, est à
\qty{11,9}{m} ; la balle doit retomber dans le carré de service, avant la ligne située à \qty{18,3}{m}.
\end{docu}
\begin{docu}{2}{Des services à 180 km/h}
Les meilleurs joueurs servent à plus de 180~km/h, soit \qty{50}{m.s^{-1}}. On néglige ici les frottements et l'effet
donné à la balle.
\end{docu}
\begin{questions}
\item La balle est frappée horizontalement à \qty{45}{m.s^{-1}}. Établir les équations horaires de son mouvement.
\item Passe-t-elle au-dessus du filet ? Où retombe-t-elle ? Conclure.
\item Quelle vitesse horizontale maximale permet à une balle frappée horizontalement de tomber dans le carré ?
\item \initiative Pour un service à \qty{50}{m.s^{-1}}, la balle doit être frappée vers le bas. Montrer qu'un angle de
$6{,}7^\circ$ sous l'horizontale convient, et vérifier qu'elle franchit le filet.
\end{questions}
\end{exercice}
\end{blocLivre}
\begin{blocLivre}{ex-15-13}
\begin{exercice}[3]{L'imprimante à jet d'encre}
\begin{docu}{1}{Des gouttes déviées}
Dans certaines imprimantes industrielles, une buse projette des gouttes d'encre de masse $m=\qty{1,0e-10}{kg}$ à
$v_0=\qty{20}{m.s^{-1}}$. Chaque goutte reçoit une charge électrique $q$ réglable, puis passe entre deux plaques de
longueur $L=\qty{1,0}{cm}$ où règne un champ uniforme $E=\qty{1,0e6}{V.m^{-1}}$, perpendiculaire à $\vec{v}_0$. Le papier est
situé à $D=\qty{1,0}{cm}$ après les plaques.
\end{docu}
\schema{15-ex13}
\begin{donnees}
Pour la goutte étudiée, $|q|=\qty{1,5e-13}{C}$ ; $g=\qty{9,81}{m.s^{-2}}$.
\end{donnees}
\begin{questions}
\item Comparer la force électrique au poids de la goutte. Quelle approximation peut-on faire ?
\item Calculer l'accélération de la goutte entre les plaques et la durée de la traversée.
\item Calculer la déviation à la sortie des plaques et la coordonnée verticale de la vitesse à cet instant.
\item \initiative Après les plaques, la goutte n'est plus soumise à aucune force notable. Calculer la déviation totale
sur le papier. Comment l'imprimante place-t-elle les gouttes à des hauteurs différentes ?
\end{questions}
\end{exercice}
\end{blocLivre}
\begin{blocLivre}{ex-15-14}
\begin{exercice}[3]{Le saut en longueur}
\begin{docu}{1}{L'impulsion}
Au moment de l'impulsion, le centre de masse d'un sauteur est à \qty{1,10}{m} du sol ; sa vitesse vaut
\qty{9,5}{m.s^{-1}}, inclinée de $22^\circ$ sur l'horizontale. À la réception, les jambes tendues vers l'avant, son centre de
masse est à \qty{0,50}{m} du sol.
\end{docu}
\begin{docu}{2}{Un angle faible}
Les sauteurs décollent sous un angle de 20 à 25°, très inférieur à 45° : en s'élevant davantage, ils perdent beaucoup de
vitesse horizontale au moment de l'impulsion.
\end{docu}
\begin{questions}
\item Établir les équations horaires du mouvement du centre de masse pendant le saut, frottements négligés.
\item Calculer l'altitude maximale atteinte par le centre de masse.
\item \initiative Calculer la durée du saut et la distance horizontale parcourue par le centre de masse.
\item Calculer la vitesse du centre de masse à la réception, à l'aide de la conservation de l'énergie mécanique.
\end{questions}
\end{exercice}
\end{blocLivre}
\begin{blocLivre}{ex-15-15}
\begin{exercice}[3]{La gouttelette de Millikan}
\begin{docu}{1}{Peser une charge électrique}
En 1909, Robert Millikan immobilise de fines gouttelettes d'huile électrisées entre deux plaques horizontales : en
ajustant la tension, il compense exactement le poids de chaque gouttelette par la force électrique.
\end{docu}
\begin{docu}{2}{Une gouttelette immobile}
Une gouttelette sphérique de rayon $r=\qty{0,80}{\micro m}$ et de masse volumique \qty{900}{kg.m^{-3}} est immobile entre
deux plaques distantes de $d=\qty{1,00}{cm}$, la plaque supérieure étant positive, pour $U=\qty{394}{V}$.
\end{docu}
\begin{docu}{3}{D'autres gouttelettes}
Charges mesurées sur quatre autres gouttelettes, en valeur absolue : \qty{3,20e-19}{C} ; \qty{4,81e-19}{C} ;
\qty{8,00e-19}{C} ; \qty{6,41e-19}{C}.
\end{docu}
\begin{donnees}
Volume d'une sphère : $\frac43\pi r^3$ ; $g=\qty{9,81}{N.kg^{-1}}$ ; $e=\qty{1,602e-19}{C}$.
\end{donnees}
\begin{questions}
\item Calculer la masse de la gouttelette, puis son poids.
\item Caractériser le champ électrique entre les plaques. Quel est le signe de la charge de la gouttelette ?
\item \initiative Déterminer la valeur absolue de la charge de la gouttelette.
\item Que suggèrent l'ensemble des valeurs obtenues ? Quelle grandeur Millikan a-t-il ainsi mesurée ?
\end{questions}
\end{exercice}
\end{blocLivre}
